% Options for packages loaded elsewhere
% Options for packages loaded elsewhere
\PassOptionsToPackage{unicode}{hyperref}
\PassOptionsToPackage{hyphens}{url}
\PassOptionsToPackage{dvipsnames,svgnames,x11names}{xcolor}
%
\documentclass[
  11pt,
  letterpaper,
  DIV=11,
  numbers=noendperiod]{scrartcl}
\usepackage{xcolor}
\usepackage[margin=1in]{geometry}
\usepackage{amsmath,amssymb}
\setcounter{secnumdepth}{5}
\usepackage{iftex}
\ifPDFTeX
  \usepackage[T1]{fontenc}
  \usepackage[utf8]{inputenc}
  \usepackage{textcomp} % provide euro and other symbols
\else % if luatex or xetex
  \usepackage{unicode-math} % this also loads fontspec
  \defaultfontfeatures{Scale=MatchLowercase}
  \defaultfontfeatures[\rmfamily]{Ligatures=TeX,Scale=1}
\fi
\usepackage{lmodern}
\ifPDFTeX\else
  % xetex/luatex font selection
\fi
% Use upquote if available, for straight quotes in verbatim environments
\IfFileExists{upquote.sty}{\usepackage{upquote}}{}
\IfFileExists{microtype.sty}{% use microtype if available
  \usepackage[]{microtype}
  \UseMicrotypeSet[protrusion]{basicmath} % disable protrusion for tt fonts
}{}
\makeatletter
\@ifundefined{KOMAClassName}{% if non-KOMA class
  \IfFileExists{parskip.sty}{%
    \usepackage{parskip}
  }{% else
    \setlength{\parindent}{0pt}
    \setlength{\parskip}{6pt plus 2pt minus 1pt}}
}{% if KOMA class
  \KOMAoptions{parskip=half}}
\makeatother
% Make \paragraph and \subparagraph free-standing
\makeatletter
\ifx\paragraph\undefined\else
  \let\oldparagraph\paragraph
  \renewcommand{\paragraph}{
    \@ifstar
      \xxxParagraphStar
      \xxxParagraphNoStar
  }
  \newcommand{\xxxParagraphStar}[1]{\oldparagraph*{#1}\mbox{}}
  \newcommand{\xxxParagraphNoStar}[1]{\oldparagraph{#1}\mbox{}}
\fi
\ifx\subparagraph\undefined\else
  \let\oldsubparagraph\subparagraph
  \renewcommand{\subparagraph}{
    \@ifstar
      \xxxSubParagraphStar
      \xxxSubParagraphNoStar
  }
  \newcommand{\xxxSubParagraphStar}[1]{\oldsubparagraph*{#1}\mbox{}}
  \newcommand{\xxxSubParagraphNoStar}[1]{\oldsubparagraph{#1}\mbox{}}
\fi
\makeatother


\usepackage{longtable,booktabs,array}
\usepackage{calc} % for calculating minipage widths
% Correct order of tables after \paragraph or \subparagraph
\usepackage{etoolbox}
\makeatletter
\patchcmd\longtable{\par}{\if@noskipsec\mbox{}\fi\par}{}{}
\makeatother
% Allow footnotes in longtable head/foot
\IfFileExists{footnotehyper.sty}{\usepackage{footnotehyper}}{\usepackage{footnote}}
\makesavenoteenv{longtable}
\usepackage{graphicx}
\makeatletter
\newsavebox\pandoc@box
\newcommand*\pandocbounded[1]{% scales image to fit in text height/width
  \sbox\pandoc@box{#1}%
  \Gscale@div\@tempa{\textheight}{\dimexpr\ht\pandoc@box+\dp\pandoc@box\relax}%
  \Gscale@div\@tempb{\linewidth}{\wd\pandoc@box}%
  \ifdim\@tempb\p@<\@tempa\p@\let\@tempa\@tempb\fi% select the smaller of both
  \ifdim\@tempa\p@<\p@\scalebox{\@tempa}{\usebox\pandoc@box}%
  \else\usebox{\pandoc@box}%
  \fi%
}
% Set default figure placement to htbp
\def\fps@figure{htbp}
\makeatother





\setlength{\emergencystretch}{3em} % prevent overfull lines

\providecommand{\tightlist}{%
  \setlength{\itemsep}{0pt}\setlength{\parskip}{0pt}}



 


\KOMAoption{captions}{tableheading}
\makeatletter
\@ifpackageloaded{caption}{}{\usepackage{caption}}
\AtBeginDocument{%
\ifdefined\contentsname
  \renewcommand*\contentsname{Table of contents}
\else
  \newcommand\contentsname{Table of contents}
\fi
\ifdefined\listfigurename
  \renewcommand*\listfigurename{List of Figures}
\else
  \newcommand\listfigurename{List of Figures}
\fi
\ifdefined\listtablename
  \renewcommand*\listtablename{List of Tables}
\else
  \newcommand\listtablename{List of Tables}
\fi
\ifdefined\figurename
  \renewcommand*\figurename{Figure}
\else
  \newcommand\figurename{Figure}
\fi
\ifdefined\tablename
  \renewcommand*\tablename{Table}
\else
  \newcommand\tablename{Table}
\fi
}
\@ifpackageloaded{float}{}{\usepackage{float}}
\floatstyle{ruled}
\@ifundefined{c@chapter}{\newfloat{codelisting}{h}{lop}}{\newfloat{codelisting}{h}{lop}[chapter]}
\floatname{codelisting}{Listing}
\newcommand*\listoflistings{\listof{codelisting}{List of Listings}}
\makeatother
\makeatletter
\makeatother
\makeatletter
\@ifpackageloaded{caption}{}{\usepackage{caption}}
\@ifpackageloaded{subcaption}{}{\usepackage{subcaption}}
\makeatother
\usepackage{bookmark}
\IfFileExists{xurl.sty}{\usepackage{xurl}}{} % add URL line breaks if available
\urlstyle{same}
\hypersetup{
  pdftitle={Relative Value Analysis of U.S. Treasury Yields and SOFR-Based Rates},
  pdfauthor={Chase Leibowitz; Dhruv Trivedi; Jose Maldonado; Felipe Sanchez-Noguera},
  colorlinks=true,
  linkcolor={blue},
  filecolor={Maroon},
  citecolor={Blue},
  urlcolor={Blue},
  pdfcreator={LaTeX via pandoc}}


\title{Relative Value Analysis of U.S. Treasury Yields and SOFR-Based
Rates}
\author{Chase Leibowitz \and Dhruv Trivedi \and Jose
Maldonado \and Felipe Sanchez-Noguera}
\date{2026-03-19}
\begin{document}
\maketitle

\renewcommand*\contentsname{Table of contents}
{
\hypersetup{linkcolor=}
\setcounter{tocdepth}{3}
\tableofcontents
}

\section{Introduction}\label{introduction}

The relationship between U.S. Treasury yields and money-market funding
rates is central to fixed-income relative value analysis. When
short-term funding rates diverge from Treasury yields of comparable
maturity, the resulting spreads can signal dislocations in credit
markets, shifts in monetary policy expectations, or changes in the
supply-demand balance for safe assets.

Traditional relative-value analysis compares Treasury yields to interest
rate swap rates (the ``swap spread''). However, swap rate data at
granular frequencies is not freely available. In this project, we
construct an analogous measure using the Secured Overnight Financing
Rate (SOFR) and its term averages as proxies for the floating-rate leg
of a swap. SOFR averages over 30-, 90-, and 180-day windows approximate
the fixed-rate equivalent of receiving SOFR over those horizons,
providing a publicly available, daily-frequency alternative to swap
rates.

We investigate three research questions:

\begin{itemize}
\tightlist
\item
  \textbf{RQ1:} Do SOFR-Treasury spreads exhibit systematic patterns
  after extreme deviations from recent norms? Specifically, do spreads
  mean-revert after breaching \(\pm2\) standard deviation thresholds, or
  do they exhibit trend continuation?
\item
  \textbf{RQ2:} How sensitive are these findings to modeling choices,
  particularly the length of the rolling normalization window used to
  compute z-scores?
\item
  \textbf{RQ3:} Does spread behavior differ across monetary policy
  regimes, specifically periods of rising versus falling or stable
  federal funds rates?
\end{itemize}

Note: These research questions have changed slightly from the original
project proposal due to complications in accessing Bloomberg Terminal
data. We hope to access this data for the final project submission and
incorporate our original research questions, but this pivot was required
at the moment due to the lack of granular data. \# Data Sources

\subsection{FRED API}\label{fred-api}

All data was sourced from the Federal Reserve Economic Data (FRED) API,
which provides free access to macroeconomic and financial time series
published by the Federal Reserve. We used the following series over the
period January 2020 through December 2025:

\begin{longtable}[]{@{}
  >{\raggedright\arraybackslash}p{(\linewidth - 4\tabcolsep) * \real{0.1928}}
  >{\raggedright\arraybackslash}p{(\linewidth - 4\tabcolsep) * \real{0.4337}}
  >{\raggedright\arraybackslash}p{(\linewidth - 4\tabcolsep) * \real{0.3735}}@{}}
\caption{FRED series used in this
analysis.}\label{tbl-series}\tabularnewline
\toprule\noalign{}
\begin{minipage}[b]{\linewidth}\raggedright
Series ID
\end{minipage} & \begin{minipage}[b]{\linewidth}\raggedright
Description
\end{minipage} & \begin{minipage}[b]{\linewidth}\raggedright
Role in Analysis
\end{minipage} \\
\midrule\noalign{}
\endfirsthead
\toprule\noalign{}
\begin{minipage}[b]{\linewidth}\raggedright
Series ID
\end{minipage} & \begin{minipage}[b]{\linewidth}\raggedright
Description
\end{minipage} & \begin{minipage}[b]{\linewidth}\raggedright
Role in Analysis
\end{minipage} \\
\midrule\noalign{}
\endhead
\bottomrule\noalign{}
\endlastfoot
DGS2 & 2-Year Treasury Constant Maturity & Short-term yield benchmark \\
DGS5 & 5-Year Treasury Constant Maturity & Medium-term yield
benchmark \\
DGS10 & 10-Year Treasury Constant Maturity & Long-term yield
benchmark \\
SOFR & SOFR Overnight Rate & Spot funding rate \\
SOFR30DAYAVG & SOFR 30-Day Average & Short-term SOFR proxy \\
SOFR90DAYAVG & SOFR 90-Day Average & Medium-term SOFR proxy \\
SOFR180DAYAVG & SOFR 180-Day Average & Long-term SOFR proxy \\
EFFR & Effective Federal Funds Rate & Regime classification \\
\end{longtable}

\subsection{Data Cleaning and
Construction}\label{data-cleaning-and-construction}

The raw data required several processing steps before analysis:

\begin{enumerate}
\def\labelenumi{\arabic{enumi}.}
\tightlist
\item
  \textbf{Type coercion.} FRED returns missing observations as the
  string \texttt{"."}. All series were coerced to numeric, converting
  these to \texttt{NaN}.
\item
  \textbf{Outer join.} Treasury and SOFR series were merged on date
  using an outer join to preserve all available observations from both
  sources.
\item
  \textbf{Forward fill.} Missing values (due to weekends, holidays, or
  publication lags) were forward-filled with a 5-day limit to avoid
  propagating stale data across extended gaps.
\item
  \textbf{Spread construction.} We defined three maturity-aligned
  spreads:

  \begin{itemize}
  \tightlist
  \item
    \texttt{spread\_2y} = SOFR30DAYAVG \(-\) DGS2
  \item
    \texttt{spread\_5y} = SOFR90DAYAVG \(-\) DGS5
  \item
    \texttt{spread\_10y} = SOFR180DAYAVG \(-\) DGS10
  \end{itemize}
\item
  \textbf{Z-score normalization.} For each spread, we computed a rolling
  z-score using a 60-day window:
  \[z_t = \frac{s_t - \bar{s}_{t,60}}{\sigma_{s,t,60}}\] where \(s_t\)
  is the spread, \(\bar{s}_{t,60}\) is the trailing 60-day rolling mean,
  and \(\sigma_{s,t,60}\) is the trailing 60-day rolling standard
  deviation.
\end{enumerate}

The final dataset contains 1,500+ daily observations across 7 rate
series, 3 spreads, and 3 z-scores.

\section{Modules}\label{modules}

\subsection{Visualization}\label{visualization}

We used \texttt{matplotlib} and \texttt{seaborn} to produce all figures
in the project. Visualization served two purposes: (1) exploratory
analysis of spread distributions and time-series behavior, and (2)
presentation of research question results. Key techniques included
time-series line plots with threshold highlighting, histograms for
distributional analysis, scatter plots for bivariate relationships, and
annotated heatmaps for correlation structure. Visualization was used
throughout the project, from initial data exploration through final
results.

\textbf{Concepts used:} distributional analysis (histograms, summary
statistics), time-series visualization, correlation heatmaps, grouped
bar charts for comparative analysis.

\subsection{Data Wrangling}\label{data-wrangling}

The core of the pipeline involved merging heterogeneous time series from
FRED into a single analysis-ready dataset. This required date alignment
across series with different publication schedules, handling of missing
values through forward-filling with safeguards, and construction of
derived features (spreads, rolling statistics, z-scores). We used
\texttt{pandas} for all data manipulation, including rolling window
calculations, conditional filtering, and reshaping for analysis. Data
wrangling was the first stage of the project and underpins all
subsequent analysis.

\textbf{Concepts used:} time-series joins, forward-fill imputation,
rolling window aggregation, feature engineering from raw financial data.

\subsection{Statistical Inference}\label{statistical-inference}

To answer our research questions, we employed one-sample t-tests to
determine whether the mean 5-day forward spread change was statistically
different from zero, conditional on z-score threshold breaches. This
allowed us to assess whether extreme z-score readings carry predictive
information about subsequent spread movements. We used
\texttt{scipy.stats.ttest\_1samp} and evaluated results at the 1\%, 5\%,
and 10\% significance levels. Statistical inference was applied in the
final analysis stage to formalize the patterns observed in exploratory
analysis.

\textbf{Concepts used:} hypothesis testing (one-sample t-test),
conditional expectation, significance levels and p-values, robustness
testing across parameter choices.

\section{Preliminary Results}\label{preliminary-results}

\subsection{Exploratory Data Analysis}\label{exploratory-data-analysis}

\subsubsection{Spread Distributions}\label{spread-distributions}

Figure~\ref{fig-histograms} shows the distribution of each spread over
the sample period. The 2-year spread is approximately centered at zero
(mean = \(-0.006\) pp), consistent with short-term SOFR and Treasury
rates tracking each other closely. The 5-year and 10-year spreads have
progressively more negative means (\(-0.043\) and \(-0.283\) pp,
respectively), reflecting the fact that longer-tenor SOFR averages are
backward-looking and tend to lag Treasury yields, which price in forward
expectations. All three distributions exhibit wide dispersion (ranges
exceeding 4 percentage points), indicating substantial variation in the
SOFR-Treasury relationship over this period.

\begin{figure}

\centering{

\includegraphics[width=1\linewidth,height=\textheight,keepaspectratio]{outputs/figures/spread_histograms.png}

}

\caption{\label{fig-histograms}Distribution of SOFR-Treasury spreads
across the three maturity pairs. Red dashed lines indicate sample
means.}

\end{figure}%

\subsubsection{Spread Time Series}\label{spread-time-series}

Figure~\ref{fig-timeseries} reveals that the spreads are far from
stationary. All three display large, persistent swings driven by the
Federal Reserve's rapid rate-hiking cycle in 2022--2023 and subsequent
stabilization. The 2-year spread plunged to approximately \(-2.5\) pp in
late 2022 as the Fed raised the funds rate faster than short-term SOFR
averages could adjust, then recovered through 2023--2024. The 5-year and
10-year spreads show similar patterns with even wider troughs, reaching
\(-3\) pp.~These regime-driven dynamics motivate the rolling z-score
approach: a static threshold on spread levels would be meaningless given
the non-stationarity, but normalizing against recent history isolates
local deviations.

\begin{figure}

\centering{

\includegraphics[width=0.85\linewidth,height=\textheight,keepaspectratio]{outputs/figures/spread_time_series.png}

}

\caption{\label{fig-timeseries}Time series of SOFR-Treasury spreads
(2020--2025). Large swings in 2022--2023 reflect the Fed's rate-hiking
cycle.}

\end{figure}%

\subsubsection{Z-Score Signals}\label{z-score-signals}

Figure~\ref{fig-zscores} shows the rolling z-scores with \(\pm2\)
threshold bands. The z-scores successfully remove the trending behavior
visible in the raw spreads, oscillating around zero with frequent but
transient breaches of the \(\pm2\) thresholds. Notably, the 2-year
z-score spent an extended period below \(-2\) in mid-2022, while the
5-year and 10-year z-scores show more frequent, shorter-duration
threshold breaches. The shaded regions highlight the signal events used
in our mean-reversion tests.

\begin{figure}

\centering{

\includegraphics[width=0.85\linewidth,height=\textheight,keepaspectratio]{outputs/figures/z_score_time_series.png}

}

\caption{\label{fig-zscores}Rolling z-scores (60-day window) for each
spread. Red shading marks \(z > +2\) episodes; green shading marks
\(z < -2\) episodes.}

\end{figure}%

\subsubsection{Correlation Structure}\label{correlation-structure}

Figure~\ref{fig-corr} shows the correlation matrix across all variables.
Treasury yields are highly correlated with each other (\(\rho > 0.96\))
and with SOFR rates (\(\rho > 0.87\)), confirming that the rate series
share a dominant common factor (the level of interest rates). The
spreads, by contrast, show moderate mutual correlation
(\(\rho \approx 0.69\)--\(0.98\)) and low correlation with rate levels
(\(\rho \approx 0.25\)--\(0.48\)), indicating that the spreads capture
information distinct from the overall rate environment.

\begin{figure}

\centering{

\includegraphics[width=0.75\linewidth,height=\textheight,keepaspectratio]{outputs/figures/correlation_matrix.png}

}

\caption{\label{fig-corr}Correlation matrix across yields, SOFR rates,
and spreads.}

\end{figure}%

\subsection{RQ1: Do Spreads Mean-Revert After Extreme
Z-Scores?}\label{rq1-do-spreads-mean-revert-after-extreme-z-scores}

We tested mean reversion by computing the 5-business-day forward change
in spread level, conditional on the z-score breaching \(\pm2\). Under
the mean-reversion hypothesis, we would expect negative forward changes
after \(z > +2\) (spread narrows back) and positive forward changes
after \(z < -2\) (spread widens back). Table~\ref{tbl-rq1} presents the
results.

\begin{longtable}[]{@{}llrrrr@{}}
\caption{RQ1: Mean 5-day forward spread change conditional on z-score
thresholds.}\label{tbl-rq1}\tabularnewline
\toprule\noalign{}
Maturity & Signal & Mean Forward (pp) & t-stat & p-value & n \\
\midrule\noalign{}
\endfirsthead
\toprule\noalign{}
Maturity & Signal & Mean Forward (pp) & t-stat & p-value & n \\
\midrule\noalign{}
\endhead
\bottomrule\noalign{}
\endlastfoot
2y & z \textless{} -2 & -0.0093 & -1.04 & 0.300 & 123 \\
2y & z \textgreater{} +2 & +0.0210 & +1.47 & 0.145 & 91 \\
5y & z \textless{} -2 & -0.0229 & -2.64 & 0.009 & 152 \\
5y & z \textgreater{} +2 & +0.0281 & +1.93 & 0.057 & 96 \\
10y & z \textless{} -2 & -0.0016 & -0.20 & 0.845 & 168 \\
10y & z \textgreater{} +2 & +0.0471 & +3.20 & 0.002 & 88 \\
\end{longtable}

\textbf{Finding: Spreads do not mean-revert --- they exhibit trend
continuation.} The signs are the opposite of what mean reversion would
predict. After \(z > +2\) (spread is wide relative to recent history),
the spread continues to widen over the next 5 days, with the effect
strongest for the 10-year maturity (+4.7 bps, \(p = 0.002\)). After
\(z < -2\), the spread continues to tighten, though this is only
statistically significant for the 5-year maturity (\(-2.3\) bps,
\(p = 0.009\)). Figure~\ref{fig-rq1} visualizes this pattern.

This result has a clear economic interpretation. The SOFR-Treasury
spread is not a pure arbitrage relationship --- it reflects structural
differences in how SOFR averages (backward-looking) and Treasury yields
(forward-looking) respond to monetary policy changes. When the Fed is
actively shifting rates, the spread trends persistently in one direction
as SOFR averages gradually catch up to where Treasury yields have
already moved. Extreme z-scores therefore signal that a repricing is
underway, not that it has gone too far.

\begin{figure}

\centering{

\includegraphics[width=0.8\linewidth,height=\textheight,keepaspectratio]{outputs/figures/rq1_mean_reversion.png}

}

\caption{\label{fig-rq1}Mean 5-day forward spread change by signal and
maturity. Positive bars after \(z > +2\) and negative bars after
\(z < -2\) indicate trend continuation, not mean reversion.}

\end{figure}%

\subsection{RQ2: Sensitivity to Normalization
Window}\label{rq2-sensitivity-to-normalization-window}

We repeated the RQ1 analysis using rolling windows of 30, 60, and 120
business days. Table~\ref{tbl-rq2} summarizes the results, and
Figure~\ref{fig-rq2} plots the relationship.

\begin{longtable}[]{@{}lrlrrr@{}}
\caption{RQ2: Selected results across normalization windows (30, 60, 120
days).}\label{tbl-rq2}\tabularnewline
\toprule\noalign{}
Maturity & Window & Signal & Mean Forward (pp) & p-value & n \\
\midrule\noalign{}
\endfirsthead
\toprule\noalign{}
Maturity & Window & Signal & Mean Forward (pp) & p-value & n \\
\midrule\noalign{}
\endhead
\bottomrule\noalign{}
\endlastfoot
2y & 30 & z \textless{} -2 & -0.0120 & 0.286 & 106 \\
2y & 60 & z \textless{} -2 & -0.0093 & 0.300 & 123 \\
2y & 120 & z \textless{} -2 & -0.0417 & \textless0.001 & 147 \\
5y & 30 & z \textless{} -2 & -0.0037 & 0.793 & 98 \\
5y & 60 & z \textless{} -2 & -0.0229 & 0.009 & 152 \\
5y & 120 & z \textless{} -2 & -0.0198 & 0.022 & 154 \\
10y & 30 & z \textgreater{} +2 & +0.0374 & 0.041 & 79 \\
10y & 60 & z \textgreater{} +2 & +0.0471 & 0.002 & 88 \\
10y & 120 & z \textgreater{} +2 & +0.0478 & \textless0.001 & 93 \\
\end{longtable}

\textbf{Finding: The trend-continuation effect is robust and strengthens
with longer windows.} The sign of the forward spread change is
consistent across all three window lengths --- no sign flips occur.
Longer windows (120 days) produce uniformly higher statistical
significance, with every maturity-signal pair reaching at least
\(p < 0.05\). This makes sense: a 120-day window is slower to
incorporate recent trend information into the rolling mean, so a
\(z > +2\) reading more reliably indicates a genuine trend rather than
short-term noise. The 30-day window, by contrast, adapts quickly and
produces weaker signals because it normalizes away much of the trend.

\begin{figure}

\centering{

\includegraphics[width=1\linewidth,height=\textheight,keepaspectratio]{outputs/figures/rq2_window_sensitivity.png}

}

\caption{\label{fig-rq2}Mean forward spread change as a function of
normalization window length. Longer windows produce larger and more
significant effects, with consistent sign across all maturities.}

\end{figure}%

\subsection{RQ3: Regime Dependence}\label{rq3-regime-dependence}

We classified each observation into a ``rising rate'' regime (EFFR
increased over the trailing 60 days) or a ``falling/stable'' regime,
then repeated the mean-reversion test within each regime.
Table~\ref{tbl-rq3} presents selected results.

\begin{longtable}[]{@{}
  >{\raggedright\arraybackslash}p{(\linewidth - 10\tabcolsep) * \real{0.1449}}
  >{\raggedright\arraybackslash}p{(\linewidth - 10\tabcolsep) * \real{0.2464}}
  >{\raggedright\arraybackslash}p{(\linewidth - 10\tabcolsep) * \real{0.1304}}
  >{\raggedleft\arraybackslash}p{(\linewidth - 10\tabcolsep) * \real{0.2899}}
  >{\raggedleft\arraybackslash}p{(\linewidth - 10\tabcolsep) * \real{0.1304}}
  >{\raggedleft\arraybackslash}p{(\linewidth - 10\tabcolsep) * \real{0.0580}}@{}}
\caption{RQ3: Mean reversion test results by monetary policy
regime.}\label{tbl-rq3}\tabularnewline
\toprule\noalign{}
\begin{minipage}[b]{\linewidth}\raggedright
Maturity
\end{minipage} & \begin{minipage}[b]{\linewidth}\raggedright
Regime
\end{minipage} & \begin{minipage}[b]{\linewidth}\raggedright
Signal
\end{minipage} & \begin{minipage}[b]{\linewidth}\raggedleft
Mean Forward (pp)
\end{minipage} & \begin{minipage}[b]{\linewidth}\raggedleft
p-value
\end{minipage} & \begin{minipage}[b]{\linewidth}\raggedleft
n
\end{minipage} \\
\midrule\noalign{}
\endfirsthead
\toprule\noalign{}
\begin{minipage}[b]{\linewidth}\raggedright
Maturity
\end{minipage} & \begin{minipage}[b]{\linewidth}\raggedright
Regime
\end{minipage} & \begin{minipage}[b]{\linewidth}\raggedright
Signal
\end{minipage} & \begin{minipage}[b]{\linewidth}\raggedleft
Mean Forward (pp)
\end{minipage} & \begin{minipage}[b]{\linewidth}\raggedleft
p-value
\end{minipage} & \begin{minipage}[b]{\linewidth}\raggedleft
n
\end{minipage} \\
\midrule\noalign{}
\endhead
\bottomrule\noalign{}
\endlastfoot
5y & rising & z \textgreater{} +2 & +0.0621 & 0.004 & 52 \\
10y & rising & z \textgreater{} +2 & +0.0904 & \textless0.001 & 53 \\
5y & falling/stable & z \textgreater{} +2 & -0.0085 & 0.649 & 44 \\
10y & falling/stable & z \textgreater{} +2 & -0.0080 & 0.732 & 35 \\
10y & falling/stable & z \textless{} -2 & +0.0603 & 0.109 & 116 \\
2y & rising & z \textless{} -2 & -0.0315 & 0.649 & 24 \\
\end{longtable}

\textbf{Finding: Trend continuation is concentrated in the rising-rate
regime; it disappears or weakly reverses in falling/stable periods.}
During rising-rate environments, the \(z > +2\) signal for the 5-year
and 10-year maturities produces large, statistically significant trend
continuation (+6.2 and +9.0 bps per 5-day period, respectively). This
aligns with the economic mechanism described in RQ1: when the Fed is
actively hiking, SOFR averages lag behind Treasury yields, creating
persistent spread widening that an extreme z-score correctly identifies
as ongoing.

In falling or stable rate regimes, the effect vanishes. The \(z > +2\)
forward changes for 5-year and 10-year spreads are slightly negative and
statistically insignificant, suggesting that when the Fed is not
actively shifting rates, extreme z-scores are indeed more noise than
signal. This finding has practical implications: any spread-based
trading signal would need to be conditioned on the monetary policy
regime to be useful.

\newpage{}

\begin{figure}

\centering{

\includegraphics[width=1\linewidth,height=\textheight,keepaspectratio]{outputs/figures/rq3_regime_analysis.png}

}

\caption{\label{fig-rq3}Mean forward spread change by regime. Trend
continuation (orange bars, \(z > +2\)) is strong during rising-rate
periods and absent during falling/stable periods.}

\end{figure}%

\section{Reflection}\label{reflection}

\subsection{Successes}\label{successes}

The pipeline successfully automates the full workflow from raw API pulls
to statistical results and publication-ready figures. Using FRED for all
data (including SOFR averages as swap-rate proxies) eliminated the need
for expensive data subscriptions while still producing economically
meaningful results. The z-score normalization framework effectively
handles the non-stationarity in spreads, and the regime-conditioning
analysis revealed a nuanced finding that would not have emerged from
unconditional tests alone.

\subsection{Challenges}\label{challenges}

The most significant challenge was the maturity mismatch between SOFR
averages and Treasury yields. A 30-day SOFR average is not truly
comparable to a 2-year Treasury yield --- it is a backward-looking
average of an overnight rate, while the Treasury yield reflects 2 years
of forward rate expectations. This mismatch likely explains why our
``spreads'' behave more like a measure of SOFR lag than a true
relative-value signal. A second challenge was the small sample size for
regime-conditioned tests; splitting the data by regime and signal
threshold reduced some cells to fewer than 30 observations.

\subsection{Collaboration}\label{collaboration}

Team members divided work between data pipeline construction (fetching,
cleaning, merging) and analysis (EDA, signal construction, statistical
testing). Team members worked mostly asynchronously.

\subsection{Next Steps}\label{next-steps}

\begin{enumerate}
\def\labelenumi{\arabic{enumi}.}
\tightlist
\item
  \textbf{Better maturity matching.} Replace SOFR averages with SOFR
  swap rates or overnight index swap (OIS) rates, which are actual term
  rates and would produce true swap spreads.
\item
  \textbf{Longer horizons.} Test 10-day, 20-day, and 60-day forward
  windows to determine whether mean reversion emerges at longer horizons
  even though 5-day behavior shows trend continuation.
\item
  \textbf{Volatility conditioning.} Test whether signal quality differs
  in high- versus low-volatility environments (e.g., using VIX or
  realized spread volatility).
\item
  \textbf{Formal backtesting.} Construct a simulated trading strategy
  based on the z-score signals and evaluate risk-adjusted returns,
  accounting for transaction costs.
\end{enumerate}

\section{AI Disclosure}\label{ai-disclosure}

Artificial Intelligence Tools: ChatGPT 5.2 via ChatGPT.com Artificial
Intelligence Tools: Claude Opus 4.6 via Claude Code Methodology: ChatGPT
was used for determining EDA methodology given the dataset Methodology:
ChatGPT was used to help find alternatives for Bloomberg data that was
not accessible for the prototype Methodology: Claude Code was used for
code syntax execution Writing -- Review and Editing: Claude Code (within
Quarto) was used for assistance in writing editing




\end{document}
